\scanpage{56}
\mspage{-3.15-}
\margin{55}
suppose
\[
\begin{array}{rl}
\text{(i)}&x^{(0)}=x^{(1)}=\cdots=x^{(\beta_0)}\ne x^{(\beta_0+1)}\\[5pt]
\text{(ii)}&x^{(0)},x^{(1)},\ldots,x^{(\beta_1)}
\quad\text{all lie on exactly one line, but}\\
&x^{(\beta_1+1)}\quad\text{is not on this line.}\\[5pt]
\text{(iii)}&x^{(0)},x^{(1)},\ldots,x^{(\beta_2)}
\quad\text{all lie on exactly one plane, but}\\
&x^{(\beta_2+1)}\quad\text{is not in this plane,}\\
&\vdots\\
\text{(n)}&x^{(0)},\ldots,x^{(\beta_{n-1})}
\quad\text{all lie in exactly one hyperplane, but}\\
&x^{(\beta_{n-1}+1)}\quad\text{is not in this hyperplane.}
\end{array}
\]
Then $x^{(0)},x^{(\beta_0+1)},\ldots,x^{(\beta_{n-1}+1)}$
do not lie in a hyperplane, hence
$D_{0,\beta_0+1,\ldots,\beta_{n-1}+1}(x)\ne0$;
and whenever $D_{\alpha_0,\alpha_1,\ldots,\alpha_n}(x)\ne0$, then
\[
\alpha_0\ge0,\quad\alpha_1\ge\beta_0+1,\quad
\alpha_2\ge\beta_1+1,\ldots,\alpha_n\ge\beta_{n-1}+1.
\]
Therefore, the condition on $\ell$ is that:
\[
(*)\qquad\ell(d_{0,\beta_0+1,\ldots,\beta_{n-1}+1})\ge0.
\]
But it is easy to check that the set of all $\{\lambda\}$'s such that
$\lambda_0\le\cdots\le\lambda_m$
and $\displaystyle\sum_{i=0}^m\lambda_i=0$ is a convex set spanned by the particular set of
$\{\lambda\}$'s given in the lemma. Therefore $(*)$ holds for some $\ell$ if and only
if it holds for one of the extreme sets of $\{\lambda_i\}$.
\uline{QED}

Removing the assumption $\lambda_0\le\cdots\le\lambda_m$, and computing $\ell$ we
conclude that the existence of the P's is equivalent to:
